\documentclass[10pt,a4paper]{amsart}
\usepackage{amsmath,amssymb,amsthm,graphicx}
\usepackage{verbatim}
\newenvironment{reference}{\comment}{\endcomment}
%\newenvironment{reference}{}{}
\newenvironment{slogan}{\comment}{\endcomment}
\newenvironment{history}{\comment}{\endcomment}

% For commutative diagrams we use Xy-pic
\usepackage[all]{xy}

% We use 2cell for 2-commutative diagrams.
\xyoption{2cell}
% Theorem environments.
%
\theoremstyle{plain}
\newtheorem{theorem}[subsection]{Theorem}
\newtheorem{proposition}[subsection]{Proposition}
\newtheorem{lemma}[subsection]{Lemma}

\theoremstyle{definition}
\newtheorem{definition}[subsection]{Definition}
\newtheorem{example}[subsection]{Example}
\newtheorem{exercise}[subsection]{Exercise}
\newtheorem{situation}[subsection]{Situation}

\theoremstyle{remark}
\newtheorem{remark}[subsection]{Remark}
\newtheorem{remarks}[subsection]{Remarks}

\numberwithin{equation}{subsection}

% Macros
%
\def\lim{\mathop{\mathrm{lim}}\nolimits}
\def\colim{\mathop{\mathrm{colim}}\nolimits}
\def\Spec{\mathop{\mathrm{Spec}}}
\def\Hom{\mathop{\mathrm{Hom}}\nolimits}
\def\Ext{\mathop{\mathrm{Ext}}\nolimits}
\def\SheafHom{\mathop{\mathcal{H}\!\mathit{om}}\nolimits}
\def\SheafExt{\mathop{\mathcal{E}\!\mathit{xt}}\nolimits}
\def\Sch{\mathit{Sch}}
\def\Mor{\mathop{\mathrm{Mor}}\nolimits}
\def\Ob{\mathop{\mathrm{Ob}}\nolimits}
\def\Sh{\mathop{\mathit{Sh}}\nolimits}
\def\NL{\mathop{N\!L}\nolimits}
\def\CH{\mathop{\mathrm{CH}}\nolimits}
\def\proetale{{pro\text{-}\acute{e}tale}}
\def\etale{{\acute{e}tale}}
\def\QCoh{\mathit{QCoh}}
\def\Ker{\mathop{\mathrm{Ker}}}
\def\Im{\mathop{\mathrm{Im}}}
\def\Coker{\mathop{\mathrm{Coker}}}
\def\Coim{\mathop{\mathrm{Coim}}}

% Boxtimes
%
\DeclareMathSymbol{\boxtimes}{\mathbin}{AMSa}{"02}

%
% Macros for moduli stacks/spaces
%
\def\QCohstack{\mathcal{QC}\!\mathit{oh}}
\def\Cohstack{\mathcal{C}\!\mathit{oh}}
\def\Spacesstack{\mathcal{S}\!\mathit{paces}}
\def\Quotfunctor{\mathrm{Quot}}
\def\Hilbfunctor{\mathrm{Hilb}}
\def\Curvesstack{\mathcal{C}\!\mathit{urves}}
\def\Polarizedstack{\mathcal{P}\!\mathit{olarized}}
\def\Complexesstack{\mathcal{C}\!\mathit{omplexes}}
% \Pic is the operator that assigns to X its picard group, usage \Pic(X)
% \Picardstack_{X/B} denotes the Picard stack of X over B
% \Picardfunctor_{X/B} denotes the Picard functor of X over B
\def\Pic{\mathop{\mathrm{Pic}}\nolimits}
\def\Picardstack{\mathcal{P}\!\mathit{ic}}
\def\Picardfunctor{\mathrm{Pic}}
\def\Deformationcategory{\mathcal{D}\!\mathit{ef}}
\usepackage{cleveref}
\crefname{originallemma}{original lemma}{original lemmas}
\Crefname{originallemma}{Original lemma}{Original lemmas}
\makeatletter
\newlabel{lemma-sheaf-with-exact-support-ext}{{0A77}{source}{}{sourcepointer-1}{}}
\newlabel{lemma-sheaf-with-exact-support-ext@cref}{{[originallemma][0][]0A77}{[1][0][]source}{}{}{}}
\makeatother
\title{The internal Hom from the structure sheaf of an effective Cartier divisor is the conormal sheaf shifted by one}
\author{The Stacks Project Authors}
\date{stacks-project:duality.tex:local-snapshot-20261012; GNU FDL 1.2 or later}
\begin{document}
\maketitle
\noindent\textit{Editorial note.} All delivered mathematical text is the author's own, copied from the pinned Stacks Project chapter file duality.tex; no statement or proof step was rewritten, reordered or removed.\par\medskip
\noindent\textit{Editorial note.} Preamble provenance: the pinned chapter file opens with an input of the shared preamble and carries no preamble text of its own, so the project's single Stacks-wide preamble staged at sources/stacks-duality/preamble.tex is excerpted verbatim under the recipe fixed by LANE\_BRIEF 22.1, as three spans, lines 8-13 (verbatim package and the reference/slogan/history environments), lines 14-18 (Xy-pic setup) and lines 154-217 (the native theorem declarations and mathematical macros). Nothing was added, removed or redefined; the wrapper class and the amsmath/amssymb/amsthm/graphicx base packages are supplied by the assembler.\par\medskip
\noindent\textit{Editorial note.} The counted statement and proof contain no cross reference, no citation and no included file, so the delivered body is self contained.\par\medskip
\noindent\textit{Editorial note.} Canonical location of the counted proposition: Stacks tag 0B4B (sources/stacks-duality/pinned-original-source.tex line 3368).\par\medskip
\noindent\textit{Allowed standard prerequisites (editorial).}\begin{itemize}
\item The sheaf with exact support and its pushforward description. Same chapter, Lemma labelled lemma-sheaf-with-exact-support-ext, cited at proof lines 3384-3388. Identifies the pushed forward internal Hom with a Hom on X.
\item The ideal sheaf sequence of an effective Cartier divisor. The short exact sequence 0 -> I -> O\_X -> i\_*O\_D -> 0 used at proof lines 3388-3392. Reduces the computation to Hom out of O\_X and out of the ideal sheaf.
\item The conormal sheaf as the restriction of the invertible sheaf of the divisor. The identities N = i\textasciicircum{}*O\_X(D) and D = zero scheme of the canonical section, used at proof lines 3402-3406. Gives the final short exact sequence and the identification with N.
\end{itemize}\medskip
\ifdefined\hypertarget\hypertarget{sourcepointer-1}{}\fi
\noindent\textit{Source reference.} Original lemma 0A77: sources/stacks-duality/pinned-original-source.tex, source lines 2339--2339. Same-chapter pointer into the author environment lemma at source line 2339; invoked at proof line 3384 of the counted proof. It is retained only as a navigation pointer; every load-bearing step is inside the delivered span.\par\medskip
\begin{lemma}
\label{lemma-compute-for-effective-Cartier}
As above, let $X$ be a scheme and let $D \subset X$ be an
effective Cartier divisor. There is a canonical isomorphism
$R\SheafHom(\mathcal{O}_D, \mathcal{O}_X) = \mathcal{N}[-1]$
in $D(\mathcal{O}_D)$.
\end{lemma}

\begin{proof}
Equivalently, we are saying that $R\SheafHom(\mathcal{O}_D, \mathcal{O}_X)$
has a unique nonzero cohomology sheaf in degree $1$ and that this
sheaf is isomorphic to $\mathcal{N}$. Since $i_*$ is exact and fully
faithful, it suffices to prove that
$i_*R\SheafHom(\mathcal{O}_D, \mathcal{O}_X)$ is isomorphic
to $i_*\mathcal{N}[-1]$. We have
$i_*R\SheafHom(\mathcal{O}_D, \mathcal{O}_X) =
R\SheafHom_{\mathcal{O}_X}(i_*\mathcal{O}_D, \mathcal{O}_X)$
by Lemma \ref{lemma-sheaf-with-exact-support-ext}. We have a resolution
$$
0 \to \mathcal{I} \to \mathcal{O}_X \to i_*\mathcal{O}_D \to 0
$$
where $\mathcal{I}$ is the ideal sheaf of $D$
which we can use to compute. Since
$R\SheafHom_{\mathcal{O}_X}(\mathcal{O}_X, \mathcal{O}_X) = \mathcal{O}_X$ and
$R\SheafHom_{\mathcal{O}_X}(\mathcal{I}, \mathcal{O}_X) = \mathcal{O}_X(D)$ by
a local computation, we see that
$$
R\SheafHom_{\mathcal{O}_X}(i_*\mathcal{O}_D, \mathcal{O}_X) =
(\mathcal{O}_X \to \mathcal{O}_X(D))
$$
where on the right hand side we have $\mathcal{O}_X$ in degree $0$
and $\mathcal{O}_X(D)$ in degree $1$. The result follows from the
short exact sequence
$$
0 \to \mathcal{O}_X \to \mathcal{O}_X(D) \to i_*\mathcal{N} \to 0
$$
coming from the fact that $D$ is the zero scheme of the canonical section
of $\mathcal{O}_X(D)$ and from the fact that
$\mathcal{N} = i^*\mathcal{O}_X(D)$.
\end{proof}

\end{document}
